\documentclass[10pt,a4paper]{amsart}
\usepackage[margin=19mm]{geometry}
\usepackage{amsmath,amssymb,mathtools}
\usepackage[scr=rsfs,cal=euler]{mathalfa}
\usepackage{xfrac}
\usepackage[unicode,hidelinks,bookmarks=false]{hyperref}
\newcommand{\ie}{i.e.\ }
\newcommand{\cf}{cf.\ }
\newcommand{\resp}{respectively\ }
\newcommand{\Subsection}{\subsection}
\providecommand{\nobreakdash}{\nobreak\hbox{-}}
\setlength{\emergencystretch}{2em}
\multlinegap0pt
\usepackage{bm}
\usepackage{bbm}
\usepackage{dsfont}
\usepackage{xspace}
\usepackage{cancel}
\usepackage{mathtools}
\usepackage{amsthm}
\makeatletter
\@ifundefined{claim}{\newtheorem{claim}{Claim}}{}
\@ifundefined{lemma}{\newtheorem{lemma}{Lemma}}{}
\makeatother
\makeatletter
\expandafter\ifx\csname mathitem\endcsname\relax
\newenvironment{mathitem}
  {\begin{list}{{$(\roman{mathitem})$}}{
   \setcounter{mathitem}{0}
   \usecounter{mathitem}
   \setlength{\topsep}{0pt plus 2pt minus 0pt}
   \setlength{\parskip}{0pt plus 2pt minus 0pt}
   \setlength{\partopsep}{0pt plus 2pt minus 0pt}
   \setlength{\parsep}{0pt plus 2pt minus 0pt}
   \setlength{\leftmargin}{25pt}
   \setlength{\itemsep}{0pt plus 2pt minus 0pt}}}
\fi
\makeatother
\title[On the maximal anti-Ramsey problem of Burr, Erd\H{o}s, Graha]{On the maximal anti-Ramsey problem of Burr, Erd\H{o}s, Graham, and S\'{o}s for $P\_4$}
\author{Zixuan Yang}
\date{Source: arXiv:2607.05896v1 (CC BY 4.0)}
\begin{document}
\maketitle

\noindent\emph{Source and licence.} On the maximal anti-Ramsey problem of Burr, Erd\H{o}s, Graham, and S\'{o}s for $P_4$. Version arXiv:2607.05896v1 (CC BY 4.0), distributed under the Creative Commons Attribution 4.0 International licence, as recorded in the arXiv licence metadata for this exact version and shown on the abstract page https://arxiv.org/abs/2607.05896v1. Full rights notice: licenses/arxiv-2607.05896-CC-BY-4.0.txt.\par\smallskip

\noindent\textit{Editorial scope.} Counted unit: the Lemma whose source label is \texttt{lem:dense-core} in the article (source0.tex lines 221--230, source label \texttt{lem:dense-core}), delivered together with the article's own complete original proof of it (source0.tex lines 232--419, one \texttt{proof} environment of 188 lines). No mathematics is written, repaired or re-derived: statement and proof are the article's own text line for line. Required context retained uncounted: none. Every definition, display and label of those lines is retained unchanged. Omitted from this package and not redistributed: 1--220, 231--231, 420--503. Nothing inside a retained excerpt is omitted or altered: every retained body is delivered exactly as the article's lines are written, so no omission receipt of any kind accompanies this package. Citations: the retained excerpts carry no citation command at all, so no bibliography entry of the article is transcribed and no reference is redistributed. Reference adaptation: none was needed and none was made; every reference inside the delivered unit resolves inside it, so no unresolved reference or citation is produced. Printed slips kept literal: no printed word, symbol, hypothesis or number of the counted statement or of its proof was altered. Layout and class: the article's own class is \texttt{article} with packages including epsf, algorithm, epsfig, amsfonts, amsgen, amsmath, amstext, amsbsy, amsopn, amsthm, lineno, color, tikz, cite; the wrapper is the lane's pinned \texttt{amsart} editorial wrapper at 10pt on A4, which restates the theorem-like environments the retained text needs and adds the article's own macro definitions that the retained text needs (none), with extra wrapper packages (bm, bbm, dsfont, xspace, cancel, mathtools). The delivered numbering is the unit's own. Assets: no retained span contains a figure environment, a graphics-inclusion command, a PDF-page inclusion, a table or a TikZ picture; no image file is copied into this package, the package declares no asset dependency and no image file is redistributed with it. Licence: version arXiv:2607.05896v1 of this article is distributed under the Creative Commons Attribution 4.0 International licence, as recorded in the arXiv licence metadata for this exact version and shown on the abstract page https://arxiv.org/abs/2607.05896v1. A full rights notice is delivered at licenses/arxiv-2607.05896-CC-BY-4.0.txt. Characters: every retained excerpt is pure ASCII, so the delivered engine, XeLaTeX with its default Latin Modern text font, reports no missing character.
\begin{lemma}\label{lem:dense-core}
Let $G$ be an $n$-vertex graph such that its complement $H=\overline G$ satisfies $|E(H)|\le n^{3/2}$.  
Let $B=\{v\in V(G): d_H(v)>8\sqrt n\}$, $U=V(G)\setminus B$, and  $F=G[U].$
For all sufficiently large $n$, the following statements hold:
\begin{itemize}
%\item [\rm{(i)}] $|U|>3n/4$, $|E(F)|> n^2/4$, $d_{\overline {F}}(v)\le 8\sqrt n$ for every $v\in U$, and  $\sum_{v\in U} d_{\overline {F}}(v)\le 2n^{3/2}$;
\item [\rm{(i)}]  if the edges of $G$ are colored such that every copy of $P_4$ in $G$ is rainbow, then every color class restricted to $E(F)$ is an induced matching in $F$;
\item [\rm{(ii)}]  if $E(F)$ can be partitioned into $q$ induced matchings,  then $q> n^2/72$.
\end{itemize}
\end{lemma}

\begin{proof}

We first consider the size of $E(F)$.

\begin{claim}\label{c4}
$|E(F)|>\frac{n^2}{4}$.
\end{claim}

\begin{proof}[Proof of Claim \ref{c4}]
Since $|E(H)|\le n^{3/2}$, the degree sum in $H$ satisfies
\[
\sum_{v\in V(G)}d_H(v)=2|E(H)|\le 2n^{3/2}.
\]
Then by the definition of $B$, we have
\[
8\sqrt n\,|B|<\sum_{v\in B}d_H(v)\le \sum_{v\in V(G)}d_H(v)\le 2n^{3/2},
\]
and therefore $|B|<n/4$, which implies that $|U|>3n/4$. 

Notice that the complement of $F$  is $H[U]$.  So we have
\begin{align}\label{eq1}
|E(F)|=\binom{|U|}{2}-|E(H[U])|
       \ge \binom{|U|}{2}-|E(H)|.
\end{align}
Substitute the inequalities $|U|>3n/4$ and $|E(H)|\le n^{3/2}$ into inequality (\ref{eq1}) to obtain the lower bound
\[
|E(F)|>\binom{3n/4}{2}-n^{3/2}
       =\frac{9n^2}{32}-\frac{3n}{8}-n^{3/2}.
\]
Since $n^2/32-n^{3/2}-3n/8\ge 0$ for  sufficiently large $n$,
we get that 
\[
|E(F)|> \frac{8n^2}{32}=\frac{n^2}{4}.
\]
Thus, the claim holds.
\end{proof}

We now prove the statement (i) by two claims. 

\begin{claim}\label{c1}
Any two adjacent edges in $F$ cannot have the same color. 
\end{claim}

\begin{proof}[Proof of Claim \ref{c1}]
By contradiction. 
Suppose that there exist two edges $xy,xz\in E(F)$ such that the colors of $xy$ and $xz$ are the same.  
Since $y\in U$, we have
$d_H(y)\le 8\sqrt n$, and consequently
\[
d_G(y)=n-1-d_H(y)\ge n-1-8\sqrt n>2,
\]
for large enough $n$. 
Thus $y$ has a neighbor $w$ in $G$ with $w\notin\{x,z\}$.  
One can see that the three edges $wy,yx,xz$ form a copy of $P_4$ in $G$, where the four vertices $w,y,x,z$ are distinct.  
Therefore, we obtain a path $P_4$ which  is not rainbow since $yx$ and $xz$ have the same color, a contradiction.  
\end{proof}


\begin{claim}\label{c2}
For any two disjoint edges  $e_1, e_2$ in $F$ with the same color, there exist no edges of $G$ connecting an endpoint of $e_1$ to an endpoint of $e_2$. 
\end{claim}

\begin{proof}[Proof of Claim \ref{c2}]

By contradiction. Suppose that there exist two disjoint edges $ab,cd\in E(F)$  with the same color such that  
 an edge of $G$ joins one endpoint of $ab$ to one endpoint of $cd$.
Then after relabeling the endpoints, without loss of generality, we may assume that $bc\in E(G)$.  
This implies that three edges $ab,bc,cd$  form a $P_4=abcd$ in $G$, and one can see that this copy of $P_4$ has two  edges with the same colors (namely $ab$ and $cd$), which  contradicts the fact that every copy of $P_4$ in $G$ is rainbow.  Therefore, the claim holds.
\end{proof}


By Claim \ref{c1},  each color class on $F$ is a matching, and then by Claim \ref{c2}, each color class
restricted to $F$ is  an induced matching in $F$. Therefore, the statement (i) holds.



In the following, we prove the statement (ii).
Let $M_1,M_2,\ldots,M_q$ be the nonempty induced matchings in $F$.
Write $ m=\sum_{i=1}^q |M_i|$. Then $m=|E(F)|$.
%By Lemma~\ref{lem:dense-core},
%\[
%d_{\overline {G[U]}}(v)\le 8\sqrt n\quad\text{for every }v\in U,
%\qquad
%\sum_{v\in U}d_{\overline {G[U]}}(v)\le 2n^{3/2},
%\qquad
%k=|E(F)|\ge \frac{n^2}{4}.
%\]
Let 
\[
\mathcal T=\{(v,i,e): v\in U,\ e\in M_i,\ v\text{ is covered by }M_i,
\text{ and }v\text{ is not an endpoint of }e\}.
\] 

We next count  the size of $\mathcal T$ by using double-counting.
For any fixed $i\in [q]$, since $M_i$ is a matching with $|M_i|$ edges, it covers exactly $2|M_i|$
vertices.  For each covered vertex $v$, there is exactly one edge of $M_i$
incident with $v$, and hence there are $|M_i|-1$ choices for an edge
$e\in M_i$ not incident with $v$.  Therefore
\begin{equation}\label{eq:T-lower}
|\mathcal T|=\sum_{i=1}^q 2|M_i|(|M_i|-1).
\end{equation}

\begin{claim}\label{c3}
For any fixed  vertex $v\in U$,  the number of triples in
$\mathcal T$ with first coordinate $v$ is at most $\binom{d_{\overline {F}}(v)}{2}$.  
\end{claim}

\begin{proof}[Proof of Claim \ref{c3}]
Let $(v,i,e)\in\mathcal T$, and write $e=xy$.  
Since $v$ is covered by $M_i$, there is a unique edge $f\in M_i$ incident with $v$, where $1\le i\le q$.  
Note that the edge $e$ is not incident with $v$. 
So $e$ and $f$ are two distinct edges of the induced matching $M_i$, where $1\le i\le q$.  
One can see that both $x$ and $y$ belong to $N_{\overline F}(v)$.
Indeed, otherwise, if $v$ were adjacent in $F$ to $x$ or to $y$, then there exists a edge of $F$  joining an
endpoint of $f$ to an endpoint of $e$, which contradicts the fact that $M_i$ is induced.

We now show that the map $(v,i,e)\mapsto \{x,y\}$ is injective.  
Indeed, the edge $e$ belongs to exactly one member of the partition, so the pair $\{x,y\}$ determines both $e$ and its index $i$.  
Hence the number of triples with first coordinate $v$ is at most the number of two-element subsets of
$N_{\overline F}(v)$, namely $\binom{d_{\overline {F}}(v)}{2}$. This completes the proof of Claim \ref{c3}.  
\end{proof}


Summing over $v\in U$ and by Claim \ref{c3} and \eqref{eq:T-lower}, we obtain that 
\begin{equation}\label{eq:key-count}
\sum_{i=1}^q 2|M_i|(|M_i|-1)=|\mathcal T|
\le \sum_{v\in U}\binom{d_{\overline {F}}(v)}{2}.
\end{equation}
%The right-hand side is small because the complement-degree in $F$ is small. 
%By   (\ref{eq0}),  for each $v\in U$, we have
%\[
%\binom{d_{\overline {F}}(v)}{2}=\frac{d_{\overline {F}}(v)(d_{\overline {F}}(v)-1)}{2}
%\le \frac{d_{\overline {F}}(v)\cdot 8\sqrt n}{2}=4\sqrt n\,d_{\overline {F}}(v).
%\]

Now we establish an upper bound of $\sum_{v\in U}\binom{d_{\overline {F}}(v)}{2}$.
By the definition of $U$, we have $d_H(v)\le 8\sqrt n$ for any $v\in U$.
 Since $\overline {F}=H[U]$, we have 
 \begin{align}\label{eq0}
 d_{\overline {F}}(v)=d_{H[U]}(v)\le d_H(v)\le 8\sqrt n.
 \end{align}
and consequently,
\begin{align}\label{eq00}
\sum_{v\in U}d_{\overline {F}}(v)=2|E(H[U])|
\le 2|E(H)|\le 2n^{3/2}.
\end{align}
By (\ref{eq0}) and (\ref{eq00}), we have
\begin{align}\label{eq5}
\sum_{v\in U}\binom{d_{\overline {F}}(v)}{2}&=\sum_{v\in U}\frac{d_{\overline {F}}(v)(d_{\overline {F}}(v)-1)}{2}\notag\\
&\le \sum_{v\in U} \frac{d_{\overline {F}}(v)\cdot 8\sqrt n}{2}\notag\\
&= 4\sqrt n\sum_{v\in U}d_{\overline {F}}(v)\notag\\
&\le 4\sqrt n\cdot 2n^{3/2}=8n^2.
\end{align}


Combining \eqref{eq:key-count} and (\ref{eq5}), we have
\begin{equation}\label{eq:ki-ki-1}
\sum_{i=1}^q (|M_i|^2-|M_i|)=\sum_{i=1}^q |M_i|(|M_i|-1)\le 4n^2.
\end{equation}
%where
%\begin{align}\label{eq6}
%\sum_{i=1}^q  |M_i|^2
%=\sum_{i=1}^q  |M_i|+\sum_{i=1}^q  |M_i|( |M_i|-1)
%=m+\sum_{i=1}^q  |M_i|( |M_i|-1).
%\end{align}
Note that 
\[
\sum_{i=1}^q |M_i|=m=|E(F)|\le \binom n2<n^2/2.
\]
So by  \eqref{eq:ki-ki-1}, we get
\begin{equation}\label{eq:square-sum}
\sum_{i=1}^q |M_i|^2=m+\sum_{i=1}^q  |M_i|( |M_i|-1)<\frac{n^2}{2}+4n^2=\frac{9n^2}{2}.
\end{equation}



Notice that
\[
m^2=\left(\sum_{i=1}^q |M_i|\right)^2\le q\sum_{i=1}^q |M_i|^2,
\]
by the Cauchy-Schwarz inequality, and $m=|E(F)|> n^2/4$ by Claim \ref{c4}. Thus, by \eqref{eq:square-sum}, we have
\[
q\ge \frac{m^2}{\sum_{i=1}^q |M_i|^2}
>\frac{(n^2/4)^2}{9n^2/2}
=\frac{n^2}{72}.
\]
This completes the proof of Lemma \ref{lem:dense-core}.
\end{proof}

\end{document}
